% Bakalárska práca - Tomáš Mucha
% Absolvovanie individuálnej odbornej praxe
\documentclass[slovak,bachelorpractice]{diploma}

% Balíky
\usepackage[autostyle=true,czech=quotes]{csquotes}
\usepackage[backend=biber, style=iso-numeric, alldates=iso]{biblatex}
\usepackage{dcolumn}
\usepackage{subfig}
\usepackage[cpp]{diplomalst}

% === METADATA ===
\ThesisAuthor{Tomáš Mucha}
\ThesisSupervisor{Ing. Pavel Dohnálek, Ph.D.}
\CzechThesisTitle{Absolvovanie individuálnej odbornej praxe}
\EnglishThesisTitle{Individual Professional Practice in the Company}
\SubmissionYear{2026}
\ThesisAssignmentFileName{ZadaniePrace.pdf}

% === POĎAKOVANIE ===
\Acknowledgement{Rád by som poďakoval vedúcemu práce Ing. Pavlovi Dohnálkovi, Ph.D. za odborné vedenie. Ďalej ďakujem konzultantovi Bc. Janovi Mitorajovi a celému tímu spoločnosti Railsformers s.r.o. za príležitosť vykonať odbornú prax a za cenné skúsenosti.}

% === ABSTRAKT SK ===
\CzechAbstract{Táto bakalárska práca popisuje priebeh individuálnej odbornej praxe v spoločnosti Railsformers s.r.o. Práca obsahuje popis firmy a pracovného zaradenia študenta, zoznam zadaných úloh s ich časovou náročnosťou a zvolený postup riešenia. Ďalej sú popísané teoretické a praktické znalosti získané počas štúdia, ktoré boli uplatnené v praxi, a tiež znalosti, ktoré študentovi chýbali. V závere sú zhrnuté dosiahnuté výsledky a celkové zhodnotenie praxe.}

\CzechKeywords{odborná prax; Railsformers; vývoj softvéru}

% === ABSTRAKT EN ===
\EnglishAbstract{This bachelor thesis describes the course of individual professional practice at Railsformers s.r.o. The thesis contains a description of the company and the student's job position, a list of assigned tasks with their time requirements, and the chosen approach to solving them. Furthermore, the theoretical and practical knowledge gained during studies that was applied in practice is described, as well as the knowledge that the student lacked. The conclusion summarizes the achieved results and overall evaluation of the practice.}

\EnglishKeywords{professional practice; Railsformers; software development}

% === SKRATKY ===
% Doplniť podľa použitých technológií
% \AddAcronym{API}{Application Programming Interface}
% \AddAcronym{REST}{Representational State Transfer}
% \AddAcronym{SQL}{Structured Query Language}

% === LITERATÚRA ===
\addbibresource{literatura.bib}

\newcolumntype{d}[1]{D{,}{,}{#1}}

% === DOKUMENT ===
\begin{document}

\MakeTitlePages

% Zoznamy (odkomentovať až budú obrázky/tabuľky)
%\listoffigures
%\clearpage
%\listoftables
%\clearpage

% === KAPITOLY ===
\chapter{Úvod}
\label{sec:Uvod}

% Zdroje použité pre túto kapitolu:
% - writing_guide.md (štruktúra)
% - updates_structured.md Fáza 1 (motivácia)
% - thesis_structure.md (obsah)
% - project_analysis.md (technický prehľad)
% - zameranie_praxe.md (pôvodná pozícia DevOps)

%% MOTIVÁCIA PRE VÝBER PRAXE %%

Pôvodne som sa na odbornú prax vo firme Railsformers~s.r.o. hlásil na pozíciu DevOps špecialistu. Náplňou mala byť migrácia projektu do verejného cloudu Azure, príprava Docker kontajnerov a~vylepšovanie CI/CD pipeline v~GitLabe. Po úvodných konzultáciách s~mentorom sa však naskytla príležitosť pracovať na projekte, ktorý prirodzene kombinoval DevOps aspekty (Docker Compose orchestrácia, kontajnerizácia služieb, správa infraštruktúry) s~full-stack vývojom (Python backend, React frontend, návrh databázy). Táto kombinácia ma zaujala viac než čisto infraštruktúrna práca, a~preto som ponuku prijal.

Projekt zameraný na IoT sledovací systém pokrýval celé spektrum moderného softvérového vývoja -- od práce s~hardvérom a~IoT protokolmi, cez návrh databázovej architektúry a~backend development, až po tvorbu interaktívneho webového rozhrania. Ďalšou významnou motiváciou bola príležitosť pracovať s~technológiami, ktoré nemajú široké zastúpenie v~bežnej výuke, ale sú v~praxi stále častejšie využívané. Technológia LoRaWAN pre bezdrôtovú komunikáciu na dlhé vzdialenosti, protokol MQTT pre IoT messaging alebo špecializované databázové rozšírenie TimescaleDB pre časové rady predstavovali pre mňa atraktívnu výzvu na rozšírenie znalostí.

Firma Railsformers má s~týmito technológiami praktické skúsenosti -- ich infraštruktúra Hedurio, na ktorej beží okrem iného MQTT broker použitý v~tomto projekte, slúži ako základ pre rôzne IoT aplikácie. Vďaka tomu som očakával kvalitný mentoring, čo sa aj potvrdilo. Nemenej dôležitým faktorom bola možnosť pracovať na projekte s~reálnou komerčnou hodnotou. Sledovací systém pre eventy má praktické využitie v~oblasti bezpečnosti, logistiky a~organizácie akcií. Možnosť vidieť, ako teoretické znalosti zo štúdia prechádzajú do funkčného produktu, bola pre mňa silnou motiváciou pre výber práve tejto praxe.

%% KONTEXT A PROBLEMATIKA IoT %%

Internet vecí (IoT -- Internet of Things) predstavuje jedno z~najdynamickejšie rastúcich odvetví informačných technológií. Podľa odhadov spoločnosti Statista dosiahne počet pripojených IoT zariadení do roku 2030 viac než 29~miliárd kusov celosvetovo~\cite{statista_iot}. Tieto zariadenia generujú obrovské množstvo dát, ktoré vyžadujú špecifické prístupy k~ukladaniu, spracovaniu a~vizualizácii.

V~kontexte sledovacích systémov pre osoby hrá kľúčovú úlohu technológia LoRaWAN (Long Range Wide Area Network). Na rozdiel od Wi-Fi alebo Bluetooth, ktoré majú dosah v~ráde desiatok metrov, LoRaWAN umožňuje komunikáciu až na vzdialenosť približne 15~kilometrov v~optimálnych podmienkach pri zachovaní nízkej spotreby energie~\cite{lorawan_spec}. To ju robí ideálnou pre aplikácie, kde je nutné pokryť rozsiahle oblasti ako festivalové areály, skladové komplexy alebo mestské štvrte, bez potreby budovať hustú sieť prístupových bodov.

Existujúce komerčné riešenia pre GPS tracking sú často nákladné, uzavreté a~neposkytujú dostatočnú flexibilitu pre špecifické požiadavky zákazníkov. Open-source alternatívy zas vyžadujú značné úsilie pri integrácii jednotlivých komponentov. Projekt, na ktorom som pracoval, mal za cieľ vytvoriť modulárne riešenie kombinujúce výhody oboch prístupov -- vlastný vývoj s~využitím otvorených technológií pri zachovaní profesionálnej kvality a~rozšíriteľnosti.

%% CIELE PRÁCE %%

Hlavným cieľom práce bolo navrhnúť a~implementovať GPS tracking systém postavený na technológii LoRaWAN, ktorý umožní real-time sledovanie osôb vybavených GPS trackermi. Systém mal spracovávať dáta z~trackerov SenseCap T1000-B komunikujúcich cez LoRaWAN gateway, ukladať ich do databázy optimalizovanej pre časové rady a~zobrazovať ich v~interaktívnej webovej aplikácii.

Medzi čiastkové ciele patrilo zabezpečenie spoľahlivého príjmu a~parsovania dát z~trackerov, návrh škálovateľnej databázovej schémy s~využitím rozšírení PostgreSQL (TimescaleDB pre časové rady, PostGIS pre priestorové dotazy), implementácia REST~API pre historické dáta aj WebSocket komunikácie pre real-time aktualizácie, a~vytvorenie používateľsky prívetivého webového rozhrania s~mapovým podkladom. Ako nadstavba bola plánovaná funkcionalita geofencingu -- definovanie virtuálnych zón s~automatickými upozorneniami pri vstupe alebo výstupe sledovaných osôb.

%% PREHĽAD NAJDÔLEŽITEJŠÍCH ZDROJOV %%

Pri práci na projekte som čerpal predovšetkým z~oficiálnych dokumentácií jednotlivých technológií. Pre backend vývoj to bola dokumentácia frameworku Flask~\cite{flask_docs} a~knižnice Socket.IO~\cite{socketio_docs}, pre databázovú vrstvu dokumentácia PostgreSQL~\cite{postgresql_docs}, TimescaleDB~\cite{timescaledb_docs} a~PostGIS~\cite{postgis_docs}. Na frontendovej strane som využíval dokumentáciu React~\cite{react_docs} a~mapovej knižnice Leaflet~\cite{leaflet_docs}. Kľúčovým zdrojom pre pochopenie IoT protokolov boli špecifikácie LoRaWAN~\cite{lorawan_spec} a~MQTT~\cite{mqtt_spec}, spolu s~datasheetmi trackerov SenseCap~\cite{sensecap_datasheet}. Nemenej dôležitá bola konzultácia s~mentorom Bc.~Janom Mitorajom, ktorý mi poskytoval praktické rady vychádzajúce z~jeho skúseností s~podobnými projektmi vo firme Railsformers.

%% ŠTRUKTÚRA PRÁCE %%

Práca je štruktúrovaná do šiestich hlavných kapitol. Prvá kapitola predstavuje firmu Railsformers~s.r.o. a~moje pracovné zaradenie v~rámci praxe. Druhá kapitola obsahuje prehľad zadaných úloh s~odhadom časovej náročnosti. Tretia kapitola, ktorá tvorí jadro práce, podrobne popisuje postup riešenia jednotlivých úloh od konfigurácie hardvéru cez implementáciu backendu a~frontendu až po testovanie. Štvrtá kapitola reflektuje teoretické a~praktické znalosti zo štúdia, ktoré som pri práci uplatnil. Piata kapitola naopak popisuje znalosti, ktoré mi chýbali a~ktoré som si musel počas praxe doplniť. Šiesta kapitola sumarizuje dosiahnuté výsledky a~obsahuje celkové zhodnotenie praxe. Záver potom prináša súhrnnú reflexiu a~návrhy na možné budúce rozšírenia systému.

\endinput

\chapter{Popis odborného zaměření firmy a pracovního zařazení studenta}
\label{sec:Firma}

% Rozsah: 3-4 strany

\section{Popis společnosti Railsformers s.r.o.}
\label{sec:PopisSpolecnosti}

\subsection{Historie a vznik}
\label{sec:Historie}
% Vznik firmy (2010), oddělení od Skvělý.CZ s.r.o.
% TODO: Doplnit obsah

\subsection{Hlavní produkty a služby}
\label{sec:Produkty}
% Hedurio - informační systém pro bezpečnostní agentury
% Vývoj informačních systémů na míru
% Mobilní aplikace
% TODO: Doplnit obsah

\subsection{IoT a Smart řešení}
\label{sec:IoT}
% Smart cities, Smart grid, LoRaWAN technologie
% TODO: Doplnit obsah

\subsection{Aktuální zaměření}
\label{sec:AktualniZamereni}
% Současný technologický stack, tým, firemní kultura
% TODO: Doplnit obsah


\section{Pracovní zařazení studenta}
\label{sec:PracovniZarazeni}

\subsection{Pozice a náplň práce}
\label{sec:Pozice}
% DevOps / Full-stack Developer
% DevOps aspekty: Docker, kontejnerizace, MQTT
% Full-stack aspekty: Python/Flask, React, PostgreSQL
% TODO: Doplnit obsah

\subsection{Mentor a spolupráce}
\label{sec:Mentor}
% Bc. Jan Mitoraj, forma spolupráce, komunikace
% TODO: Doplnit obsah

\subsection{Projekt}
\label{sec:Projekt}
% LoRaWAN GPS Tracking System
% Účel, technologie, pracovní prostředí
% TODO: Doplnit obsah

\subsection{Přínos pozice}
\label{sec:PrinosPozice}
% Komplexní praxe, celý stack, reálný projekt
% TODO: Doplnit obsah

\endinput

\chapter{Seznam úkolů zadaných studentovi s vyjádřením časové náročnosti}
\label{sec:Ulohy}

% Rozsah: 1-2 strany

\section{Přehled úkolů}
\label{sec:PrehledUkolu}

% Hlavní tabulka úkolů s časovou náročností
% TODO: Doplnit tabulku
%
% | # | Úkol | Popis | Technologie | Hodiny |
% |---|------|-------|-------------|--------|
% | 1 | Hardware Setup | Gateway, trackery | LoRaWAN | ~15h |
% | 2 | MQTT + Parser | TLS broker, parsování | Python | ~25h |
% | 3 | Databáze | Schema, TimescaleDB, PostGIS | PostgreSQL | ~20h |
% | 4 | Backend API | Flask REST, WebSocket | Flask | ~30h |
% | 5 | Frontend | React, Leaflet, 3 módy | React | ~35h |
% | 6 | Geofencing | Polygony, alerting | PostGIS | ~15h |
% | 7 | Optimalizace | Bezpečnost, edge cases | Všechny | ~15h |
% | 8 | Testování | Simulátor, real-world | Python | ~10h |
% | CELKEM | | | | ~165h |


\section{Stručný popis úkolů}
\label{sec:PopisUkolu}

\subsection{Hardware Setup}
% Konfigurace LoRaWAN gateway a GPS trackerů
% TODO: Doplnit popis

\subsection{MQTT a Parser}
% Implementace MQTT klienta s TLS a parseru
% TODO: Doplnit popis

\subsection{Databáze}
% Návrh schématu s TimescaleDB a PostGIS
% TODO: Doplnit popis

\subsection{Backend API}
% REST API a real-time WebSocket
% TODO: Doplnit popis

\subsection{Frontend}
% React aplikace s mapou a 3 režimy
% TODO: Doplnit popis

\subsection{Geofencing}
% Virtuální zóny s detekcí a notifikacemi
% TODO: Doplnit popis

\subsection{Optimalizace a bezpečnost}
% TLS, CORS, edge cases, performance
% TODO: Doplnit popis

\subsection{Testování}
% Simulátor a testování s reálnými trackery
% TODO: Doplnit popis

\endinput

\chapter{Postup riešenia zadaných úloh}
\label{sec:Riesenie}

% TODO: Napísať detailný popis riešenia (15-25 strán)
% Podľa zadania: bod 3 - zvolený postup riešenia
% Toto je JADRO práce - najdlhšia kapitola

Táto kapitola detailne popisuje postup riešenia jednotlivých zadaných úloh.

% ============================================
\section{Úloha 1: Názov}
\label{sec:ukol1}

\subsection{Zadanie a ciele}
% Čo bolo cieľom úlohy?

\subsection{Analýza problému}
% Ako ste problém analyzovali?

\subsection{Návrh riešenia}
% Aký bol navrhnutý postup?

\subsection{Implementácia}
% Ako ste riešenie implementovali?
% Ukážky kódu, screenshoty, diagramy

% Príklad vloženia kódu:
% \begin{lstlisting}[label=src:example1, caption={Popis kódu}]
% // váš kód tu
% \end{lstlisting}

% Príklad vloženia obrázka:
% \begin{figure}[htbp]
%     \centering
%     \includegraphics[width=0.8\textwidth]{Figures/screenshot.png}
%     \caption{Popis obrázka}
%     \label{fig:example1}
% \end{figure}

\subsection{Testovanie}
% Ako ste riešenie testovali?

\subsection{Výsledky a zhodnotenie}
% Aký bol výsledok? Čo sa podarilo/nepodarilo?

% ============================================
\section{Úloha 2: Názov}
\label{sec:ukol2}

% Rovnaká štruktúra ako Úloha 1

\subsection{Zadanie a ciele}

\subsection{Analýza problému}

\subsection{Návrh riešenia}

\subsection{Implementácia}

\subsection{Testovanie}

\subsection{Výsledky a zhodnotenie}

% ============================================
% Pridať ďalšie sekcie pre ďalšie úlohy...

\endinput

\chapter{Teoretické a praktické znalosti uplatnené v praxi}
\label{sec:Znalosti}

% TODO: Napísať reflexiu znalostí (1 strana)
% Podľa zadania: bod 4

Táto kapitola popisuje teoretické a praktické znalosti získané v priebehu štúdia, ktoré študent uplatnil počas odbornej praxe.

\section{Znalosti z programovania}
% Ktoré predmety/kurzy vám pomohli?
% - Algoritmy a dátové štruktúry
% - Objektovo orientované programovanie
% - Databázové systémy
% - Webové technológie
% - ...

\section{Znalosti z matematiky a teoretických predmetov}
% - Diskrétna matematika
% - Lineárna algebra
% - ...

\section{Soft skills}
% - Tímová práca
% - Prezentačné zručnosti
% - Komunikácia
% - Time management
% - ...

\section{Zhrnutie}
Celkovo možno povedať, že štúdium na FEI VŠB-TUO poskytlo dobrý základ pre...

\endinput

\chapter{Znalosti a zručnosti chýbajúce študentovi}
\label{sec:Chybajuce}

% TODO: Napísať sebareflexiu (1 strana)
% Podľa zadania: bod 5

Táto kapitola identifikuje znalosti a zručnosti, ktoré študentovi v priebehu odbornej praxe chýbali a ktoré by bolo vhodné do študijného programu zaradiť alebo prehĺbiť.

\section{Technické znalosti}
% Čo ste museli naštudovať sami?
% - Konkrétne technológie/frameworky
% - Špecifické nástroje
% - Best practices z priemyslu

\section{Praktické zručnosti}
% - Práca v tíme na reálnom projekte
% - Code review
% - CI/CD
% - ...

\section{Odporúčania pre študijný program}
% Čo by ste odporučili zaradiť do výučby?

\endinput

\chapter{Dosažené výsledky v průběhu odborné praxe a její celkové zhodnocení}
\label{sec:Vysledky}

% Rozsah: 3-4 strany

\section{Shrnutí implementované funkcionality}
\label{sec:Funkcionalita}
% Kompletní tracking systém:
% - Real-time tracking
% - Historical playback
% - Heatmap vizualizace
% - Geofencing
% - POI management
% - Device management
% - Battery monitoring
% TODO: Doplnit obsah


\section{Metriky projektu}
\label{sec:Metriky}

% Tabulka metrik:
% | Metrika | Hodnota |
% |---------|---------|
% | Backend řádků | ~1,500 |
% | Frontend komponent | 13 |
% | REST API endpointů | 28+ |
% | WebSocket eventů | 5 |
% | Databázových tabulek | 8 |
% | Docker services | 4 |
% | Odpracovaných hodin | ~165 |
% TODO: Doplnit tabulku


\section{Technická kvalita}
\label{sec:TechnickaKvalita}
% Modulární architektura, Separation of concerns
% Škálovatelný databázový návrh
% TLS zabezpečení
% TODO: Doplnit obsah


\section{Feedback od mentora}
\label{sec:Feedback}
% Citáty mentora
% TODO: Doplnit citáty


\section{Co bych udělal jinak}
\label{sec:CoJinak}

\subsection{Dřívější real-world testování}
\label{sec:RealWorldTesting}
% TODO: Doplnit obsah

\subsection{Automatizované testy}
\label{sec:AutomatizovaneTesty}
% TODO: Doplnit obsah

\subsection{Dokumentace během vývoje}
\label{sec:DokumentaceBehem}
% TODO: Doplnit obsah

\subsection{CI/CD pipeline}
\label{sec:CICD}
% TODO: Doplnit obsah

\endinput

\chapter{Záver}
\label{sec:Zaver}

Táto práca dokumentuje priebeh mojej odbornej praxe vo firme Railsformers s.r.o. V~súlade so zadaním som postupne popísal odborné zameranie firmy, zoznam zadaných úloh, zvolený postup riešenia, uplatnené aj chýbajúce znalosti a~dosiahnuté výsledky s~celkovým zhodnotením.

Pracovné prostredie vo firme Railsformers bolo mimoriadne podporujúce. Ako malá firma s~priateľskou atmosférou mi poskytla priestor na samostatnú prácu, no zároveň som mal vždy k~dispozícii pomoc skúsenejších kolegov. Flexibilný prístup k~organizácii práce, možnosť home-office a~pravidelné konzultácie s~mentorom vytvorili ideálne podmienky pre učenie sa nových technológií. Ocenil som predovšetkým ochotu tímu vysvetľovať nielen technické aspekty, ale aj širší kontext rozhodnutí a~osvedčené postupy z~praxe.

Absolvovanie odbornej praxe mi prinieslo realistický pohľad na prácu v~softvérovej firme. Pred praxou som mal predstavu o~vývoji softvéru založenú prevažne na školských projektoch. Reálna práca však ukázala, že rovnako dôležité ako samotné programovanie sú aj komunikácia, dokumentácia, riešenie neočakávaných problémov a~schopnosť prispôsobiť sa meniacim požiadavkám. Táto skúsenosť mi dala sebavedomie vstúpiť do profesionálneho IT prostredia.

Budúcim študentom jednoznačne odporúčam absolvovať odbornú prax. Je to príležitosť overiť si teoretické znalosti v~praxi, získať cenné kontakty a~zistiť, ktorým oblastiam IT sa chcete venovať. Práca na reálnom projekte s~konkrétnymi deadlinmi a~požiadavkami je skúsenosť, ktorú žiadny školský predmet nedokáže plnohodnotne nahradiť.

\endinput


% === LITERATÚRA ===
\printbibliography[title={Literatúra}, heading=bibintoc]

% === PRÍLOHY ===
% \appendix
% \chapter{Obsah priloženého média}
\label{sec:AppendixMedia}

Priložené médium obsahuje nasledovnú štruktúru:

\begin{itemize}
    \item \texttt{/src/backend/} -- Python backend (Flask, MQTT, parser)
    \item \texttt{/src/frontend/} -- React frontend
    \item \texttt{/docker-compose.yml} -- Docker orchestrácia
    \item \texttt{/database/} -- SQL migrácie a schéma
    \item \texttt{/README.md} -- Inštrukcie na spustenie
    \item \texttt{/docs/} -- API dokumentácia
    \item \texttt{/thesis/} -- Zdrojové súbory tejto práce (LaTeX)
\end{itemize}

\endinput


\end{document}
